Web search engines index the web on a schedule that lags behind real-world change.
When a user queries a time-sensitive fact---a rental price, a software release version,
a recently updated statistic---the returned result may reflect a stale version of the
source page. This temporal mismatch is practically important yet poorly quantified:
no standard benchmark exists for measuring how often web search results are out of
date, or by how much.

We introduce \textbf{FreshState}, a benchmark construction pipeline that records
ground-truth change events on live web pages and uses them to evaluate whether
language models propagate stale information from outdated search snippets.
FreshState deploys daily monitors on candidate URLs, extracts answer-bearing
values (prices, version strings) via domain-specific extractors, and records
change events when those values shift between consecutive snapshots.
We apply FreshState to two high-churn domains---rental apartment listings
(Craigslist, SF Bay Area) and open-source software releases (GitHub)---tracking
1,628 URLs over a 14-day window.

In a controlled \emph{snippet-swap} experiment, we feed language models either
a fresh or stale search snippet and measure whether they echo the outdated value.
Across three models (GPT-4o, GPT-4o-mini, Claude Sonnet), stale context causes
models to reproduce the outdated value in 100\% of cases, compared to
100\% current-value accuracy when given fresh context.
The software domain produces change events at a rate of ${\sim}17\%$ of tracked
URLs per day, while apartment listings exhibit a ${\sim}1\%$ daily price-change
rate alongside a ${\sim}23\%$ daily expiration rate, confirming that staleness is
a pervasive, quantifiable phenomenon. We release FreshState's pipeline, extractor
library, and seed datasets to support reproducible evaluation of temporal freshness
in retrieval-augmented generation systems.
